% ===================================================================
%  TABEL Nr. 13 — Die hotel. — Die restaurant. — Die kafee.
%  Meme regime que les precedents. Tableau a UNE SEULE COLONNE, avec
%  DEUX NOTES a deux rangs differents : t13-noto-1 ouvre le fichier,
%  avant meme l'apparat, et renvoie au tableau 6 pour le detail de la
%  chambre ; t13-noto-2 se place entre le 8 et l'intertitre du
%  dejeuner. Les deux sont en gras ET en italique dans le
%  fac-simile ; on reproduit les deux.
%
%  LE 9 EST UN MENU, ET LE GRAS Y COURT SUR LES PLATS AVEC LEUR
%  PONCTUATION. Les traductions sortent la virgule du gras — le
%  neerlandais le fait — parce que la regle des colonnes est le gras
%  sur le terme seul. Le fac-simile ido, lui, ne bouge pas.
%
%  ET « tripi segun la modo di Caen » RESTE UN PLAT DE CAEN. On ne
%  cherche pas d'equivalent local : le menu est un document sur ce
%  qu'on mangeait dans cet hotel-la.
%
%  ATTENTION AU MOT « grumo » : il parait deux fois, aux renvois 19
%  et 54, et c'est le meme metier — le chasseur, le joggie. Deux
%  renvois pour un mot, ce n'est pas une faute a corriger.
% ===================================================================

%%K t13-noto-1 noto
\VUnotes{143.2mm}{%
\textit{\VUgras{(*) Sien tabel nr. 6 vir die besonderhede van die kamer.}}%
}

%%K t13-apar-1 apar
\VUsaut{3.0mm}
\VUtitre{15.0pt}{60}{TABEL Nr. 13}
\VUsaut{1.6mm}
\VUpk{10.2pt}{40}{Die hotel. --- Die restaurant.}
\VUsaut{2.0mm}
\VUpk{10.2pt}{40}{Die kafee.}
\VUsaut{3.4mm}

%%K t13-01-1 p
1. --- Nadat ek gisteraand in Royan aangekom het, het ek my intrek geneem in die
«\textit{Hotel van die Oseaan}» wat deur 'n vriend by my aanbeveel is.

%%K t13-01-2 p
'n Mens het my 'n mooi \VUgras{kamer} \textsuperscript{(2)} gegee, op die tweede
\VUgras{verdieping} \textsuperscript{(1)} waar ek baie goed geslaap het (*).

%%K t13-02-1 p
2. --- Vanoggend, toe ek my kamer uitgegaan het waarin die \VUgras{diensbode}
\textsuperscript{(3)} die ontbyt gebring het, het ek die \VUgras{rookkamer}
\textsuperscript{(5)} binnegegaan wat ook as \VUgras{leeskamer} \textsuperscript{(4)}
gebruik word om die \VUgras{koerante} \textsuperscript{(6)} te lees, terwyl ek 'n
\VUgras{sigaret} \textsuperscript{(8)} gerook het. Ná die lees het ek met die
\VUgras{hysbak} \textsuperscript{(9)} afgegaan en deur die stad gaan wandel.

%%K t13-03-1 p
3. --- Omdat ek vergeet het om die \VUgras{beampte} \textsuperscript{(12)} van die hotel
te vra op watter uur die maaltye aan die \VUgras{gemeenskaplike tafel}
\textsuperscript{(11)} opgedien word, het ek vroeg teruggekom en dit benut om in die
\VUgras{badkamer} \textsuperscript{(10)} van die hotel te gaan bad. Daarna het ek weer na
die \VUgras{kassa} \textsuperscript{(15)} gegaan om een van my vriende te bel. Maar die
\VUgras{telefoon} \textsuperscript{(13)} was besig; toe die \VUgras{kassierster}
\textsuperscript{(14)}, wat 'n \VUgras{rekening} \textsuperscript{(17)} in die
\VUgras{reisigersregister} \textsuperscript{(16)} ingeskryf het, my verleentheid sien,
het sy die \VUgras{portier} \textsuperscript{(18)} gevra om die \VUgras{joggie}
\textsuperscript{(19)} te stuur om my boodskap \VUgras{per fiets} \textsuperscript{(20)}
te doen.

%%K t13-04-1 p
4. --- Ek het haar bedank en uitgegaan om 'n fooitjie te gee aan die \VUgras{kruier}
\textsuperscript{(24)} (die sjouer) wat van die stasie af op sy \VUgras{handkar}
\textsuperscript{(25)} my \VUgras{hoededoos} \textsuperscript{(26)}, \VUgras{my koffer}
\textsuperscript{(27)} en \VUgras{my reissak} \textsuperscript{(28)} gebring het. Die
\VUgras{briewebesteller} \textsuperscript{(21)} wat pas aangekom het, het my 'n
\VUgras{brief} \textsuperscript{(22)} gegee.

%%K t13-05-1 p
5. --- Ek het nog 'n rukkie op die drumpel van die deur gebly, naby die
\VUgras{briewebus} \textsuperscript{(23)}, om na die verkeer op die straat te kyk. Die
\VUgras{kondukteur} \textsuperscript{(30)} van die \VUgras{omnibus}
\textsuperscript{(29)} het op sy wa die \VUgras{koffer} \textsuperscript{(31)} gelaai van
'n \VUgras{reisiger} \textsuperscript{(32)} van wie die \VUgras{hotelhouer}
\textsuperscript{(33)} afskeid geneem het. 'n Jong \VUgras{telegrambesteller}
\textsuperscript{(35)} het sonder haas 'n \VUgras{telegram} \textsuperscript{(34)} vir 'n
gas van die hotel gebring; 'n \VUgras{toeris} \textsuperscript{(36)} met sy
\VUgras{fototoestel} \textsuperscript{(38)} oor die skouer het sy \VUgras{reisgids}
\textsuperscript{(37)} geraadpleeg.

%%K t13-tit-1 sub
\VUsaut{4.0mm}
\VUpk{10.2pt}{40}{Etenstyd.}
\VUsaut{3.0mm}

%%K t13-06-1 p
6. --- Voor die hek het 'n sorgelose \VUgras{boodskapper} \textsuperscript{(39)}
gedommel tussen sy \VUgras{draaghaak} \textsuperscript{(40)} en sy \VUgras{poetsdoos}
\textsuperscript{(41)}, terwyl 'n jong \VUgras{wasvrou} \textsuperscript{(42)} wat 'n
\VUgras{mandjie} \textsuperscript{(43)} linnegoed gedra het, gelag het oor die plegtige
gesig van die \VUgras{tafelbeampte} \textsuperscript{(44)} wat by die deur gestaan het.

%%K t13-07-1 p
7. --- Toe ek die \VUgras{kok} \textsuperscript{(45)} sien wat sy \VUgras{kombuis}
\textsuperscript{(47)} in die \VUgras{kelderverdieping} \textsuperscript{(48)} uitgekom
het om 'n \VUgras{skottelwasser} \textsuperscript{(46)} 'n boodskap te laat doen, het ek
onthou dat die uur van die ete naby is en ek het die restaurant binnegegaan, terwyl ek
oor die binnehof gestap het waarop 'n \VUgras{motoris} \textsuperscript{(50)} sy
\VUgras{motor} \textsuperscript{(49)} weggesit het, terwyl 'n \VUgras{masjinis}
\textsuperscript{(52)} volgens die aanwysings van die \VUgras{fietsryer}
\textsuperscript{(53)} 'n \VUgras{petroldriewiel} \textsuperscript{(51)} herstel het wat
pas 'n ongeluk gehad het.

%%K t13-08-1 p
8. --- Ek het 'n jong \VUgras{joggie} \textsuperscript{(54)} gevra wat
\VUgras{bierglase} \textsuperscript{(55)} gedra het, of die gemeenskaplike tafel onder
die veranda staan, en op sy bevestigende antwoord het ek plek by een van die tafels
ingeneem en my 'n uitstekende maaltyd laat gee.

%%K t13-noto-2 noto
\VUnotes{143.2mm}{%
\textit{\VUgras{(*) Met behulp van die geregtelys wat in die woordelys opgeneem is, sal
die onderwyser die onderstaande spyskaart maklik kan varieer.}}%
}

%%K t13-tit-2 sub
\VUsaut{4.0mm}
\VUpk{10.2pt}{40}{'n Uitstekende ete.}
\VUsaut{3.0mm}

%%K t13-09-1 p
9. --- Die kelner het my die spyskaart (*) van die ete en die wynlys gebring. Ek het
gekies en as \VUgras{voorgereg garnale} met \VUgras{botter} bestel, en, as
tussengereg, \VUgras{pens op die wyse van Caen}. Ek het nie \VUgras{vis} geëet nie, maar
'n porsie \VUgras{lamsboud} gevra, en, as \VUgras{groenteskottel, spinasie} met room.
Daarna, omdat geen van die \VUgras{tussengeregte} my beval het nie, het ek verskillende
nageregte laat bring, \VUgras{kaas}, \VUgras{vrugte} en \VUgras{koekies}. Ek het boonop
'n uitstekende Saint-Émilion-\VUgras{wyn} bygevoeg, en baie tevrede oor my maaltyd het ek
die tafel verlaat nadat ek die \VUgras{kelner} \textsuperscript{(57)} 'n mooi
\VUgras{fooitjie} \textsuperscript{(59)} gegee het.

%%K t13-10-1 p
10. --- Toe die \VUgras{bestuurder} \textsuperscript{(60)} my gereed sien om te vertrek,
het hy my voorgestel om na die balkon te gaan om die pragtige panorama van die
\VUgras{Oseaan} \textsuperscript{(62)} te bewonder. In die verte het 'n mens \VUgras{vissersbootjies}
\textsuperscript{(63)}, \VUgras{seilskepe} \textsuperscript{(64)} en 'n enorme
\VUgras{pakketboot} \textsuperscript{(65)} begin sien waarvan die swart buitelyn teen die
wit \VUgras{wolke} \textsuperscript{(67)} van die \VUgras{horison} \textsuperscript{(66)}
uitstaan. 'n Ligte bries het die \VUgras{vlag} \textsuperscript{(70)} laat golf wat bo
die \VUgras{uithangbord} \textsuperscript{(69)} van die \VUgras{saal}
\textsuperscript{(68)} vasgesteek is.

%%K t13-tit-3 sub
\VUsaut{3.0mm}
\VUpk{10.2pt}{40}{Vermaak.}
\VUsaut{3.0mm}

%%K t13-11-1 p
11. --- Die skouspel was so boeiend dat ek besluit het om môre op die terras van die
kafee te klim en 'n \VUgras{toneelkyker} \textsuperscript{(71)} of 'n \VUgras{verkyker}
\textsuperscript{(72)} saam te neem.

%%K t13-12-1 p
12. --- Toe ek die balkon verlaat het, het ek na die kafee van die hotel gegaan om 'n
\VUgras{drankie} \textsuperscript{(73)} te sluk, terwyl ek 'n \VUgras{partytjie biljart}
\textsuperscript{(75)} gespeel het. Sonder om stil te hou het ek deur die kafeesaal
gegaan waarin baie verbruikers \VUgras{skaak} \textsuperscript{(78)}, \VUgras{dambord}
\textsuperscript{(79)}, \VUgras{dominostene} \textsuperscript{(80)}, \VUgras{triktrak}
\textsuperscript{(81)} of \VUgras{kaarte} \textsuperscript{(82)} gespeel het, en ek het
reguit na die \VUgras{biljartsaal} \textsuperscript{(74)} opgegaan.

%%K t13-13-1 p
13. --- Toe die partytjie klaar was, het ek na die grondverdieping afgegaan en by die
kelner 'n \VUgras{koevert} \textsuperscript{(84)}, \VUgras{papier}
\textsuperscript{(83)}, 'n \VUgras{posseël} \textsuperscript{(85)}, 'n \VUgras{pen}
\textsuperscript{(87)} en \VUgras{ink} \textsuperscript{(86)} gevra om 'n brief te skryf.

%%K t13-13-2 p
Daarna het ek 'n \VUgras{koetsier} \textsuperscript{(92)} geroep wie se \VUgras{rytuig}
\textsuperscript{(88)} voor die kafee stilgehou het. Ek het hom na die prys volgens die
ure of vir een rit gevra, en toe ons dit oor die prys eens was, het hy vir my die
\VUgras{rytuigdeur} \textsuperscript{(89)} oopgemaak en my ingehelp. Hy het weer op die
\VUgras{bok} \textsuperscript{(91)} geklim, die perde vinnig met 'n \VUgras{sweepie}
\textsuperscript{(93)} laat wegtrek en ons het vir 'n bekoorlike rit vertrek.
\VUsaut{6.0mm}
\VUfilet{22mm}
